% Standalone problem document, generated from the adjacent README.md.
% Compile with XeLaTeX. Mathematical content is maintained in Markdown.
\documentclass[11pt,a4paper]{article}
\usepackage[fontsize=10.5pt]{scrextend}
\usepackage[margin=22mm,headheight=15pt,headsep=9mm,footskip=11mm]{geometry}
\usepackage{fontspec}
\setmainfont{texgyrepagella}[Extension=.otf,UprightFont=*-regular,BoldFont=*-bold,ItalicFont=*-italic,BoldItalicFont=*-bolditalic]
\setsansfont{texgyreheros}[Extension=.otf,UprightFont=*-regular,BoldFont=*-bold,ItalicFont=*-italic,BoldItalicFont=*-bolditalic]
\setmonofont{texgyrecursor}[Extension=.otf,UprightFont=*-regular,BoldFont=*-bold,ItalicFont=*-italic,BoldItalicFont=*-bolditalic,Scale=MatchLowercase]
\usepackage{amsmath,amssymb}
\usepackage{unicode-math}
\setmathfont{texgyrepagella-math.otf}
\usepackage{xcolor,microtype,enumitem,needspace,fancyhdr}
\usepackage{bookmark}
\definecolor{ink}{HTML}{17344B}
\definecolor{accent}{HTML}{197D80}
\definecolor{muted}{HTML}{596773}
\hypersetup{colorlinks=true,linkcolor=accent,urlcolor=accent,pdftitle={MI-02 - Bourin's
crossed-Heinz
inequality},pdfauthor={Open Problems in Numerical Linear Algebra}}
\urlstyle{same}
\setlength{\parindent}{0pt}
\setlength{\parskip}{5pt plus 1pt minus 1pt}
\linespread{1.04}
\setlength{\emergencystretch}{3em}
\widowpenalty=10000
\clubpenalty=10000
\displaywidowpenalty=10000
\allowdisplaybreaks[1]
\setlist{nosep,leftmargin=1.5em}
\providecommand{\tightlist}{\setlength{\itemsep}{0pt}\setlength{\parskip}{0pt}}
\providecommand{\pandocbounded}[1]{#1}
\pagestyle{fancy}
\fancyhf{}
\fancyhead[L]{\sffamily\footnotesize\color{muted}OPEN PROBLEMS IN NUMERICAL LINEAR ALGEBRA}
\fancyhead[R]{\sffamily\bfseries\footnotesize\color{accent}MI-02}
\fancyfoot[L]{\parbox[b]{0.9\textwidth}{\sffamily\fontsize{8}{10}\selectfont\color{muted}Editorial ratings; a bounded search cannot certify that a problem remains unsolved.\\Literature check: 2026-09-10}}
\fancyfoot[R]{\sffamily\footnotesize\color{muted}\thepage}
\renewcommand{\headrulewidth}{0.3pt}
\renewcommand{\headrule}{\hbox to\headwidth{\color{accent}\leaders\hrule height \headrulewidth\hfill}}
\makeatletter
\renewcommand{\section}{\Needspace{5\baselineskip}\@startsection{section}{1}{0pt}{12pt plus 2pt minus 2pt}{4pt}{\sffamily\bfseries\large\color{ink}}}
\renewcommand{\subsection}{\Needspace{4\baselineskip}\@startsection{subsection}{2}{0pt}{9pt plus 1pt minus 1pt}{3pt}{\sffamily\bfseries\normalsize\color{ink}}}
\makeatother
\setcounter{secnumdepth}{0}
\begin{document}
{\sffamily\small\color{muted}Matrix inequalities and norms\par}
\vspace{3pt}
{\raggedright\hyphenpenalty=10000\exhyphenpenalty=10000\sffamily\bfseries\fontsize{21}{24}\selectfont\color{ink}Bourin's
crossed-Heinz inequality\par}
\vspace{7pt}
{\sffamily\small\textbf{Difficulty:} extreme\quad\textbf{Importance:} interesting
to the community\par}
{\sffamily\small\bfseries\color{accent}Status: Partially resolved\par}
\vspace{4pt}
{\color{accent}\hrule height 0.8pt}
\vspace{6pt}
\textbf{Rating rationale:} A longstanding sharp bound for all unitarily
invariant norms remains beyond the known interpolation range; it would
strengthen matrix-power and matrix-mean estimates.

\subsection{Problem statement}\label{problem-statement}

For every integer \(n\ge1\), positive definite
\(A,B\in\mathbb C^{n\times n}\), real \(t\in[0,1]\), and unitarily
invariant norm \(|||\cdot|||\), is

\[|||A^tB^{1-t}+B^tA^{1-t}|||\le|||A+B|||?\]

Powers use the spectral functional calculus. A norm is unitarily
invariant when \(|||UXV|||=|||X|||\) for all unitary \(U,V\). The
positive definite formulation avoids ambiguity at exponent zero; for
\(0<t<1\) it includes the positive semidefinite case by continuity.

\subsection{Why it matters}\label{why-it-matters}

This asks for a sharp bound on noncommuting products of fractional
matrix powers. Bounds of this kind support analysis of positive definite
matrix computations and the comparison of matrix means.

\subsection{References}\label{references}

\begin{enumerate}
\def\labelenumi{\arabic{enumi}.}
\tightlist
\item
  J.-C. Bourin, \emph{Matrix subadditivity inequalities and
  block-matrices}, International Journal of Mathematics 20(6) (2009),
  679--691; the crossed-Heinz question.
  \href{https://arxiv.org/abs/0805.1954}{Preprint},
  \href{https://doi.org/10.1142/S0129167X09005509}{DOI}.
\item
  F. Kittaneh and É. Ricard, \emph{On a question of Bourin}, Linear
  Algebra and its Applications 710 (2025), 356--362; central parameter
  interval. \href{https://doi.org/10.1016/j.laa.2025.02.004}{DOI}.
\item
  T. Zhang, \emph{Bourin-type inequalities for \(\tau\)-measurable
  operators in fully symmetric spaces}.\\
  arXiv:2602.00358v1 (30 January 2026), §1, equation (1.3), Theorem 1.3.
  \href{https://arxiv.org/html/2602.00358}{Primary text}.
\end{enumerate}

\subsection{Status check --- 2026-09-10}\label{status-check-2026-09-10}

The bound with constant one is proved for \(t\in[1/4,3/4]\); the
endpoints \(0,1\) are immediate for positive definite matrices. The
displayed assertion remains unresolved in the two outer open intervals.
The latest 2026 source gives a larger explicit constant there. Searches
included \textit{Bourin crossed Heinz conjecture 2026},
\textit{On a question of Bourin Kittaneh Ricard}, and
\textit{Bourin inequality proof counterexample 2025 2026}. No
resolution of the remaining norm inequality was found. The stronger
assertion for individual singular values has counterexamples and is not
the present problem.

\textbf{Audit update (2026-09-10):} Rechecked Zhang's Theorem 1.3 and
searched for a later crossed-Heinz proof or counterexample. Constant one
is established on \([1/4,3/4]\); the larger constant on the outer
intervals leaves the stated question there unresolved. This is a bounded
literature check, not a proof that no solution exists.
\end{document}
